\chapter{Review of Related Literature} \label{sec:rrl}
Recent advances in autonomous robotics and computer vision have explored various approaches for perception, communication, and image enhancement. In the context of underwater scene understanding, three main areas are particularly relevant: vision-language models (VLMs) for robotic perception, semantic communication methods for low-bandwidth environments, and specialized techniques for underwater image enhancement.

\section[Paull, Huang, Seto \& Leonard]{Communication-Constrained Multi-AUV Cooperative SLAM~\cite{lowbandwith}}
Multi-AUV cooperative SLAM enables collaborative localization through direct encounters (relative range measurements) and indirect encounters (shared observations of seabed features). However, underwater acoustic channels operate at approximately 1,500 m/s propagation speed with throughputs of only 10--100 bytes per second, 20--50\% packet dropout, and TDMA constraints, making raw sensor data transmission impractical. A graph-based C-SLAM algorithm addresses this by locally marginalizing vehicle pose estimates via the Schur complement, condensing accumulated observations into compact virtual measurements. A sparsification step based on constrained convex optimization further reduces packet size by minimizing the Kullback-Leibler divergence between the original and approximate distributions while enforcing a predetermined sparse Jacobian structure. The resulting payload scales linearly with observed features and remains constant in time even under dropped packets. Simulation with two AUVs demonstrates bounded localization error at 50\% dropout with a maximum packet size of 192 bytes per 10-second interval.

These severe bandwidth constraints directly motivate ATHENA's semantic captioning approach: if even compressed SLAM packets must remain under 200 bytes to be acoustically viable, transmitting compact natural language scene descriptions is far more practical than any form of image transmission.

\section[Luo, Zhou, Zeng, Jiang, Qi, Xiang, Pang \& Tang]{Robotics Perception and Control: Key Technologies and Applications~\cite{perception}}
This comprehensive review surveys nine sensor categories used in robotic control systems, including IMUs, visual sensors, LiDAR, tactile sensors, and EMG sensors, categorized as either proprioceptive (measuring internal robot states) or exteroceptive (gathering information from the external environment). Vision sensors occupy a central role, supporting object recognition, visual SLAM, autonomous navigation, medical surgery, and deep learning-based robotic grasping across industrial, medical, and exploratory domains.

The survey situates VLMs as an emerging class of exteroceptive perception module that converts visual input into structured semantic representations for downstream tasks such as navigation and scene understanding. ATHENA adopts this paradigm by deploying a VLM as an onboard sensor module that produces textual descriptions suitable for transmission over low-bandwidth acoustic channels, analogous to how onboard ATR systems convert raw sonar returns into compact feature representations.

\section[Li \& Zhang]{Real-Time Vision--Language Analysis for Autonomous Underwater Drones~\cite{li2}}
A cloud-edge framework for real-time VLM analysis in AUVs uses Qwen2.5-VL, a 72-billion-parameter model, by offloading inference to a cloud server while retaining lightweight preprocessing on the edge device. The edge performs frame sampling at 1 fps and a three-stage image enhancement pipeline (color correction, backscatter reduction, white balance), while communication relies on a 5 Mbps fiber-optic link with adaptive JPEG quality scaling and bounded frame buffering. Simulation across pipeline inspection, coral reef monitoring, and wreck investigation scenarios demonstrates consistent scene comprehension with end-to-end latencies near 1 second.

However, the reliance on cloud infrastructure and tethered optical fiber makes this approach unsuitable for untethered AUV operations, underscoring the need for a fully onboard, lightweight inference pipeline such as the one adopted in ATHENA.

\section[Adang, Low, Shorinwa \& Schwager]{SINGER: An Onboard Generalist Vision-Language Navigation Policy for Drones~\cite{singer}}
SINGER addresses open-vocabulary drone navigation by integrating a photorealistic language-embedded 3D Gaussian Splatting simulator, an RRT*-based trajectory generator, and a lightweight visuomotor policy trained on synthetic data. The system uses CLIPSeg to convert RGB images into semantic heatmaps as policy inputs, enabling generalization across unseen environments. To achieve real-time inference on the NVIDIA Jetson Orin Nano 8GB, CLIPSeg is traced to ONNX and executed via the CUDA runtime at 12 Hz. Hardware experiments demonstrate 76.67\% goal proximity success across 30 trials, outperforming a classical PD controller baseline by 23.33\% with 10\% fewer collisions.

This result establishes a key precedent for ATHENA's edge deployment strategy, confirming that VLM-based semantic understanding can run in real time on the Jetson Orin Nano without cloud offloading and validating the hardware feasibility of onboard inference in resource-constrained robotic platforms.

\section[Marafioti et al.]{SmolVLM: Redefining Small and Efficient Multimodal Models~\cite{marafioti4}}
SmolVLM is a family of compact multimodal models that systematically investigates architectural configurations, tokenization strategies, and data composition to maximize VLM performance at minimal memory footprints. The architecture pairs a SigLIP vision encoder with a SmolLM2 language backbone and uses pixel shuffle compression to reduce visual tokens by a factor of 16. Three variants are produced: SmolVLM-256M ($<$ 1 GB RAM), SmolVLM-500M (1.2 GB), and SmolVLM-2.2B (4.9 GB). The 2.2B model achieves 59.8\% average score across nine benchmarks using only 4.9 GB VRAM, outperforming Qwen2VL-2B which requires 13.7 GB. SmolVLM-256M, operating under 1 GB RAM, significantly surpasses the 80B Idefics model despite a 300-fold parameter gap.

SmolVLM-Instruct serves as the lightweight backbone evaluated in ATHENA, representing the sub-1 GB end of the model scale spectrum and confirming that competitive captioning performance is achievable within the memory budgets of embedded AUV hardware.

\section[J. Li, D. Li, Savarese \& Hoi]{BLIP-2: Bootstrapping Language-Image Pre-Training with Frozen Image Encoders and Large Language Models~\cite{blip2}}
BLIP-2 proposes a compute-efficient pre-training strategy that bootstraps vision-language learning from frozen image encoders and frozen LLMs using a lightweight Querying Transformer (Q-Former) with only 188 million trainable parameters, 54 times fewer than Flamingo80B, while surpassing it by 8.7\% on zero-shot VQAv2. Q-Former employs 32 learned query vectors to extract task-relevant visual information through cross-attention, serving as an information bottleneck between the vision and language components. The model is pre-trained in two stages: first learning visual representations through contrastive, generative, and matching objectives on a 129M-image dataset; then projecting Q-Former outputs as soft visual prompts to a frozen LLM for language generation. On zero-shot captioning, BLIP-2 achieves a CIDEr score of 121.6 on the NoCaps validation set.

However, systematic analyses identify BLIP-2 as producing formulaic outputs with low lexical diversity and limited reasoning consistency~\cite{hani2026systematic}. In ATHENA, BLIP-2 OPT 2.7B serves as the intermediate-scale backbone in the ablation study, providing a reference point between SmolVLM and Qwen3.5-VL.

\section[Hani, Tagougui \& Kherallah]{Toward Human-Level Understanding: A Systematic Review of Vision-Language Models for Image Captioning~\cite{hani2026systematic}}
A systematic review evaluates eight VLMs for image captioning under zero-shot conditions on MS-COCO, combining automatic metrics (CLIPScore, perplexity, lexical diversity) with human evaluation across 1,000 captions per model. LLaVA achieves the highest CLIPScore (0.757), while BLIP-2 exhibits the lowest lexical diversity and highest perplexity (97.1), characteristic of its conservative Q-Former bottleneck. A structured WH-question analysis across six semantic dimensions (Subject Detail, Action Detail, Location Detail, Time Detail, Manner Detail, and Purpose/Reasoning) reveals a clear progression from early models producing basic subject-action-location statements to advanced models generating multi-component narratives approaching human-level descriptive richness.

Critically, the review identifies that VLMs pre-trained on terrestrial datasets consistently struggle with domain-shifted inputs, exhibiting degraded performance on imagery with fundamentally different visual statistics. This finding directly motivates ATHENA's fine-tuning strategy: standard zero-shot VLMs cannot produce accurate underwater scene descriptions without supervised adaptation on domain-specific data.

\section[NVIDIA]{TensorRT-LLM: Embedded Inference for Large Vision-Language Models~\cite{tensorrt}}
The TensorRT Edge-LLM framework provides tooling for deploying large language and vision-language models on embedded hardware through aggressive quantization and inference optimization. Targeting devices such as the NVIDIA Jetson family, the framework applies INT4 and INT8 weight quantization, KV-cache compression, and attention kernel fusion to reduce memory footprint and inference latency. Its 4-bit weight quantization achieves approximately 4$\times$ memory reduction compared to FP16 inference, making large VLMs viable on platforms with limited unified memory.

Support for Qwen3.5-VL through this framework justifies its selection as a viable ATHENA backbone despite its larger parameter count, with 4-bit NF4 quantization via the \textit{bitsandbytes} library confirming deployment feasibility on the Jetson Orin hardware platform.

\section[H. Li, Wang, Zhang, L. Li \& Ren]{Underwater Image Captioning: Challenges, Models, and Datasets~\cite{mainref}}
Underwater image captioning is hindered by three challenges: models pre-trained on terrestrial images create a systematic domain gap that causes misidentification of marine objects; underwater images exhibit degraded quality from wavelength-dependent absorption, backscatter, and turbidity; and existing underwater datasets either lack captions or are too small for training. The proposed UICM-SOFF model addresses these by using a physics-based submergence GAN to generate 80,000 degraded training images, pre-training a Meta-ResNeXtor feature extractor via meta-learning, and fusing scene-level features with YOLOv8 object-level features through a cross-attention Transformer encoder. The constructed dataset comprises 3,176 images across ten categories (predominantly fish, coral reefs, and divers) with five human-written captions each, totaling 15,880 captions.

This dataset serves as the primary training and evaluation benchmark for ATHENA, and UICM-SOFF's scores (BLEU-1: 0.7249, BLEU-4: 0.3953, ROUGE-L: 0.5888, CIDEr: 1.2831) define the state-of-the-art baseline against which ATHENA's results are directly compared.

\section[W. Chen, Xu, H. Chen, X. Zhang, Qin, Y. Zhang \& Han]{Semantic Communication Based on Large Language Model for Underwater Image Transmission~\cite{chen1}}
An LLM-based semantic communication framework for underwater image transmission represents image content as compact semantic information rather than pixel data, reducing transmitted data to approximately 0.8\% of the original image size. The semantic encoder passes the image to a visual LLM to generate text encoding high-level content, identifies key image regions via bounding box prediction at reduced resolution, and aggressively resamples the remaining background; the combined payload is approximately 27,525 bytes versus the original 3,145,728 bytes. At the receiver, a diffusion model reconstructs the full image from three guidance signals conditioned on the recovered text and key regions. Evaluation on the SUIM dataset shows the highest CLIP scores across all tested SNR levels (0--18 dB).

Both this framework and ATHENA share the core insight that language-based representations serve as an effective acoustic communication payload to overcome the fundamental bandwidth limitations of underwater channels.

\section[Zhang et al.]{An Underwater Acoustic Semantic Communication Approach to Underwater Image Transmission~\cite{zhang3}}
This work proposes an acoustic semantic communication framework that extracts semantic features from underwater images using deep neural networks and encodes them for delivery over bandwidth-limited, high-latency acoustic channels. Semantic encoding separates task-relevant content from perceptually redundant detail, achieving significant bandwidth savings while maintaining downstream task performance under challenging acoustic conditions.

These findings reinforce the practical viability of semantic representations for AUV communication and provide quantitative support for ATHENA's design choice to transmit VLM-generated captions rather than compressed images over acoustic channels.

\section[Luo, Chen \& Guo]{Semantic Communications: Overview, Open Issues, and Future Research Directions~\cite{luo2022semantic}}
A broad survey of semantic communication systems identifies three communication levels: technical (reliable bit transmission), semantic (accurate meaning transfer), and effectiveness (achieving desired receiver effects). Semantic systems encode only semantically relevant content while discarding redundant data, which is particularly valuable in bandwidth-constrained environments such as underwater acoustic channels where available bitrates are orders of magnitude lower than those needed for raw image transmission.

This framework underpins the rationale for ATHENA's use of VLM-generated captions as the communication payload: by transmitting meaning rather than pixel representations, ATHENA operates as a semantic communication system at the second level, achieving efficiency gains through language-level scene representation.

\section[Shi, Xiao, Li \& Xie]{From Semantic Communication to Semantic-Aware Networking: Model, Architecture, and Open Problems~\cite{shi2021semantic}}
Semantic-aware networking decouples semantic encoding from physical-layer transmission, enabling network resources to be allocated based on semantic importance rather than treating all bits uniformly. The proposed architecture introduces a semantic source encoder, a channel encoder, and a semantic decoder that reconstructs intended meaning without necessarily recovering original data exactly. Unlike conventional compression, which minimizes bitrate subject to distortion constraints, semantic encoding minimizes the information needed to achieve a specified level of semantic fidelity.

This decoupled architecture supports ATHENA's goal of generating compact captions that can be prioritized for transmission over low-bandwidth acoustic channels, where operator situational awareness depends on scene content rather than image fidelity.

\section[Xie, Qin, Li \& Juang]{Deep Learning Enabled Semantic Communication Systems~\cite{xie2021deep}}
End-to-end deep learning methods can jointly optimize semantic encoding and channel coding, outperforming conventional bit-based communication under low SNR conditions. The system trains a semantic encoder-decoder pair with the channel as a non-trainable intermediate, learning which input aspects are most important for downstream tasks and encoding them more robustly. Experiments show significant improvements in semantic fidelity at low SNR values typical of underwater acoustic channels.

This lends support to ATHENA's use of text captions as a robust underwater transmission format, since natural language descriptions constitute a compact semantic representation that preserves task-relevant content more reliably than compressed image data under noisy channel conditions.

\section[Huang, Tao, Gao \& Lu]{Deep Learning-Based Image Semantic Coding for Semantic Communications~\cite{huang2021deep}}
A deep learning-based framework transmits semantic image representations rather than compressed pixel data through joint source-channel coding. The semantic encoder learns scene-level abstractions appropriate for downstream tasks while discarding irrelevant texture detail, and the decoder uses generative modeling to reconstruct images from the compact semantic code. The approach scales to very low bitrate regimes corresponding to the available bandwidth of underwater acoustic channels.

While this framework transmits learned continuous-valued semantic features, ATHENA extends the same principle to discrete natural language representations: VLM-generated captions encode the semantically relevant aspects of underwater scenes in a format that is both human-interpretable and efficiently transmissible over acoustic channels.

\section[J. Wu, C. Wu, Lin, Yoshinaga, Zhong, Chen \& Ji]{Semantic Segmentation-Based Semantic Communication System for Image Transmission~\cite{wu2023semantic}}
This system applies semantic segmentation as the transmission representation for image communication, extracting pixel-level semantic labels instead of raw pixels. Segmentation maps require significantly lower bitrates while preserving object category distribution and scene structure, sufficient for many monitoring and inspection tasks. The system demonstrates robust performance under channel noise conditions typical of wireless and acoustic links.

This approach contextualizes ATHENA's choice of natural language captions as a compact alternative: while segmentation maps encode spatial structure, captions encode semantic content in a directly human-readable form that can be interpreted without specialized decoding equipment, making them well-suited for transmission to surface operators.

\section[Akkaynak \& Treibitz]{Sea-thru: A Method For Removing Water From Underwater Images~\cite{seathru}}
Underwater images suffer from systematic color distortion and contrast loss caused by wavelength-dependent absorption and scattering. Sea-thru addresses these by proposing a physics-based image formation model that separates signal attenuation from backscatter as functions of range-to-scene depth. By estimating and subtracting the backscatter contribution, then compensating for channel-specific attenuation, the method recovers true object reflectances validated against photographs taken in air.

However, the requirement for depth information, typically obtained from multi-image structure-from-motion reconstruction, limits real-time onboard applicability. The method also does not address contrast or sharpness degradation beyond color correction. These limitations motivate the more comprehensive multi-step enhancement pipeline adopted in ATHENA, which addresses color, contrast, and sharpness without requiring depth information.

\section[Wang, Sun \& Ren]{Meta Underwater Camera: A Smart Protocol for Underwater Image Enhancement~\cite{7steppipeline}}
Underwater images suffer from three categories of degradation: color deviations from wavelength-dependent attenuation, low contrast from light scattering, and detail blur from backscatter. The Meta Underwater Camera (MUC) protocol addresses all three through a cascaded seven-step pipeline: (1) attenuated channel compensation, (2) white balance via YCbCr mean method, (3) tone mapping, (4) HSL saturation adjustment, (5) contrast stretching, (6) gamma correction, and (7) high-pass fusion. A reinforcement learning agent trained with PPO dynamically configures the eight pipeline parameters using a 512-dimensional multi-color-space feature vector and a reward based on UCIQE and UIQM quality metrics. Comparative evaluation against nine methods on the UIEB dataset demonstrates superior image quality across diverse scene types.

This seven-step protocol forms the backbone of ATHENA's preprocessing module, with one modification: the contrast stretching step is replaced with CLAHE for improved local contrast preservation in heterogeneous underwater scenes.

\section[Lopez-Vazquez, Lopez-Guede, Chatzievangelou et al.]{Deep Learning Based Deep-Sea Automatic Image Enhancement and Animal Species Classification~\cite{3step}}
A combined image enhancement and classification pipeline targets deep-sea benthic megafauna from footage captured at 870 meters depth. The enhancement proceeds in two phases: classical preprocessing (white balance, gamma correction, CLAHE) followed by a 14-layer residual CNN that generates enhanced images. The full pipeline achieves the highest image quality scores (UIQM of 2.585, UCIQE of 2.406), and classification accuracy on CNN-enhanced images reaches 79.44\% compared to 66.44\% on unenhanced data.

The substantial improvement in downstream task accuracy confirms that preprocessing enhancement is a necessary precondition for effective visual understanding, directly supporting the inclusion of a dedicated image enhancement pipeline in ATHENA's architecture.

\section[Du, Sun, K. Wang, Yang \& C. Wang]{Underwater Image Enhancement Method Based on Entropy Weight Fusion~\cite{3step1}}
An entropy weight fusion framework constructs three complementary enhancement branches: gray-world white balance with adaptive gamma correction, CLAHE with bilateral filtering, and Single-Scale Retinex with guided filtering in HSV space. The three outputs are fused using information-theoretic weights derived from UIQM, entropy, and average gradient, requiring no reference images. Evaluation on the UIEBD dataset demonstrates superior performance over competing methods across greenish, bluish, hazy, and limited-illumination scenes.

These results reinforce ATHENA's multi-step enhancement design: no single technique can address the full spectrum of underwater degradation, and combining complementary approaches achieves robust improvement across the diverse conditions encountered in real AUV deployments.

\section[Wong, Yu, Ho \& Paramesran]{Comparative Analysis of Underwater Image Enhancement Methods in Different Color Spaces~\cite{clahe}}
A comparative analysis evaluates standard contrast stretching and CLAHE as representative global and local contrast enhancement techniques for underwater images. CLAHE consistently outperforms contrast stretching by operating on local tiles independently, enhancing dark regions without over-brightening well-illuminated areas, while its contrast limit parameter prevents excessive noise amplification.

These findings justify the substitution of CLAHE for the original contrast stretching step in ATHENA's preprocessing pipeline, as its local adaptivity is particularly beneficial for underwater images that frequently contain both dark shadowed regions and brightly illuminated foreground subjects.
